% --- CORRECCION: Forzar interpretacion raw de la entrada ---
\UseRawInputEncoding
\documentclass[aspectratio=169, 10pt, xcolor=table]{beamer}

% --- Configuracion del Tema ---
\usetheme{Madrid}
\usecolortheme{dolphin}

% --- Paquetes necesarios ---
\usepackage[utf8]{inputenc}
\usepackage[spanish]{babel}
\usepackage{amsmath, amssymb}
\usepackage{booktabs}
\usepackage{graphicx}
\usepackage{listings}
\usepackage{tikz}
\usetikzlibrary{arrows.meta, positioning, shapes.geometric}
\usepackage{etoolbox}

% --- Ruta de Imagenes ---
\graphicspath{{./imagenes/}}

% --- Estilo para codigo Python (no se usa mucho aqui, pero por si acaso) ---
\definecolor{codegreen}{rgb}{0,0.6,0}
\definecolor{codegray}{rgb}{0.5,0.5,0.5}
\definecolor{codepurple}{rgb}{0.58,0,0.82}
\definecolor{backcolour}{rgb}{0.95,0.95,0.92}
\lstset{
    language=SQL,
    backgroundcolor=\color{backcolour},
    commentstyle=\color{codegreen},
    keywordstyle=\color{blue},
    numberstyle=\tiny\color{codegray},
    stringstyle=\color{codepurple},
    basicstyle=\ttfamily\scriptsize,
    breaklines=true,
    frame=single,
    showstringspaces=false,
    literate={á}{{\'a}}1 {é}{{\'e}}1 {í}{{\'i}}1 {ó}{{\'o}}1 {ú}{{\'u}}1 {ñ}{{\~n}}1 {ü}{{\"u}}1
}

% --- Pie de pagina personalizado ---
\makeatletter
\setbeamertemplate{footline}{
  \leavevmode%
  \hbox{%
  \begin{beamercolorbox}[wd=.333333\paperwidth,ht=2.25ex,dp=1ex,center]{author in head/foot}%
    \usebeamerfont{author in head/foot}\insertshortauthor
  \end{beamercolorbox}%
  \begin{beamercolorbox}[wd=.333333\paperwidth,ht=2.25ex,dp=1ex,center]{title in head/foot}%
    \usebeamerfont{title in head/foot}Sesion 1
  \end{beamercolorbox}%
  \begin{beamercolorbox}[wd=.333333\paperwidth,ht=2.25ex,dp=1ex,right]{date in head/foot}%
    \usebeamerfont{date in head/foot}BBDD \hfill \insertframenumber/\inserttotalframenumber \hspace*{0.3cm}
  \end{beamercolorbox}}%
  \vskip0pt%
}
\makeatother

% --- Datos de la portada ---
\title[SESION 1 BBDD]{\textbf{Sesion 1}}
\subtitle{El Universo del Dato y la Infraestructura Física}
\author[Prof. C. Sanchez]{Carlos Cesar Sanchez Coronel}
\institute{\textbf{BASES DE DATOS}}
\date{2026}

% --- Eliminar barra de navegacion ---
\setbeamertemplate{navigation symbols}{}

% --- Indice al inicio de cada seccion ---
\AtBeginSection[]{
  \begin{frame}
    \frametitle{Contenido}
    \tableofcontents[currentsection]
  \end{frame}
}

\begin{document}

% --- Portada ---
\begin{frame}
    \titlepage
\end{frame}

% --- Indice general ---
\begin{frame}{Contenido}
    \tableofcontents
\end{frame}

% ============================================================================
% SECCION 1: DATO, INFORMACION Y CONOCIMIENTO
% ============================================================================
\section{Dato, Informacion y Conocimiento}

\begin{frame}{Jerarquia del dato}
    \begin{columns}[t]
        \column{0.33\textwidth}
        \textbf{Dato}
        \begin{itemize}
            \item Unidad minima, hecho objetivo.
            \item Ej: $1010$, ``Perez'', \texttt{2025-03-15}, \texttt{true}.
        \end{itemize}
        \column{0.33\textwidth}
        \textbf{Informacion}
        \begin{itemize}
            \item Dato contextualizado.
            \item Ej: $1010$ es ID del cliente Juan Perez.
        \end{itemize}
        \column{0.33\textwidth}
        \textbf{Conocimiento}
        \begin{itemize}
            \item Informacion asimilada para accion.
            \item Ej: Clientes con ID similar compran en marzo $\Rightarrow$ campana.
        \end{itemize}
    \end{columns}
    \vspace{0.3cm}
    \centering
    \begin{tikzpicture}[scale=0.8]
        \draw[fill=blue!20] (0,0) rectangle (2,1) node[midway] {Dato};
        \draw[fill=green!20] (3,0) rectangle (5,1) node[midway] {Informacion};
        \draw[fill=red!20] (6,0) rectangle (8,1) node[midway] {Conocimiento};
        \draw[->, thick] (2,0.5) -- (3,0.5);
        \draw[->, thick] (5,0.5) -- (6,0.5);
    \end{tikzpicture}
\end{frame}

% ============================================================================
% SECCION 2: CLASIFICACION DE DATOS
% ============================================================================
\section{Clasificacion de Datos}

\begin{frame}{Tipos segun su estructura}
    \begin{table}
        \centering
        \begin{tabular}{l|l}
            \toprule
            Tipo & Ejemplo \\
            \midrule
            \textbf{Estructurados} & Tabla Clientes (ID, Nombre, Fecha) \\
            \textbf{Semiestructurados} & JSON, XML, YAML \\
            \textbf{No estructurados} & Imagenes, audio, video, texto libre \\
            \bottomrule
        \end{tabular}
    \end{table}
    \begin{itemize}
        \item Datos no estructurados requieren procesamiento especial (embeddings, CNN, transformers).
    \end{itemize}
\end{frame}

\begin{frame}{Tipos de datos elementales}
    \begin{columns}[t]
        \column{0.5\textwidth}
        \begin{itemize}
            \item \textbf{Numericos}: INT, BIGINT, FLOAT, DECIMAL
            \item \textbf{Texto}: CHAR, VARCHAR, TEXT
            \item \textbf{Fecha/hora}: DATE, TIME, TIMESTAMP
        \end{itemize}
        \column{0.5\textwidth}
        \begin{itemize}
            \item \textbf{Binarios}: BLOB (imagenes, audio)
            \item \textbf{Booleanos}: BOOL
        \end{itemize}
    \end{columns}
\end{frame}

% ============================================================================
% SECCION 3: CODIFICACION DE CARACTERES
% ============================================================================
\section{Codificacion de Caracteres}

\begin{frame}{UTF-8: el estandar universal}
    \begin{itemize}
        \item Longitud variable (1-4 bytes), cubre todo Unicode.
        \item Compatible con ASCII (1 byte para ingles).
        \item Soporta todos los idiomas y emojis.
        \item Resuelve limitaciones de ASCII y Latin-1.
    \end{itemize}
    \begin{exampleblock}{Ejemplo}
        Caracter ``\~n'': en UTF-8 se codifica como 2 bytes (C3 B1).
    \end{exampleblock}
\end{frame}

% ============================================================================
% SECCION 4: FORMATOS DE ALMACENAMIENTO
% ============================================================================
\section{Formatos de Almacenamiento}

\begin{frame}{Comparacion de formatos}
    \begin{table}
        \centering
        \begin{tabular}{l|c|c|c}
            \toprule
            Formato & Legible & Eficiencia analitica & Soporte esquema \\
            \midrule
            CSV & Si & Baja & No \\
            JSON & Si & Media & Implicito \\
            Avro & No & Media & Si (evolutivo) \\
            Parquet & No & Alta & Si \\
            \bottomrule
        \end{tabular}
    \end{table}
    \begin{itemize}
        \item \textbf{Parquet}: ideal para data lakes (columnar, compresion).
    \end{itemize}
\end{frame}

% ============================================================================
% SECCION 5: METADATA
% ============================================================================
\section{Metadata}

\begin{frame}{Los datos que hablan de los datos}
    \begin{block}{Tipos de metadata}
        \begin{description}
            \item[Descriptiva:] descubrimiento (titulo, autor).
            \item[Estructural:] estructura (formato, tipo de dato).
            \item[Administrativa:] gestion (fecha creacion, permisos).
        \end{description}
    \end{block}
    \begin{itemize}
        \item \textbf{Data lineage}: rastreo de origen y transformaciones, basado en metadata.
        \item Ayuda a auditar, depurar y entender el impacto de cambios.
    \end{itemize}
\end{frame}

% ============================================================================
% SECCION 6: LOGS Y DATOS TEMPORALES
% ============================================================================
\section{Logs y Datos Temporales}

\begin{frame}{Registros de eventos y datos efimeros}
    \begin{columns}[t]
        \column{0.5\textwidth}
        \textbf{Logs}
        \begin{itemize}
            \item Registros secuenciales de eventos.
            \item Usos: auditoria, recuperacion, rendimiento.
            \item Ej: binlog de MySQL.
        \end{itemize}
        \column{0.5\textwidth}
        \textbf{Datos temporales}
        \begin{itemize}
            \item Caches, tablas temporales.
            \item Ciclo de vida corto.
            \item Ej: \texttt{CREATE TEMPORARY TABLE} en SQL.
        \end{itemize}
    \end{columns}
\end{frame}

% ============================================================================
% SECCION 7: HARDWARE Y JERARQUIA DE MEMORIA
% ============================================================================
\section{Hardware y Jerarquia de Memoria}

\begin{frame}{Jerarquia de memoria}
    \begin{table}
        \centering
        \begin{tabular}{l|c}
            \toprule
            Nivel & Latencia aproximada \\
            \midrule
            Cache L1 & 1 ns \\
            RAM (buffer pool) & 100 ns \\
            NVMe & 20-50 $\mu$s (20.000-50.000 ns) \\
            SSD & 0.1-0.2 ms (100.000-200.000 ns) \\
            HDD & 4-10 ms (4.000.000-10.000.000 ns) \\
            \bottomrule
        \end{tabular}
    \end{table}
    \begin{itemize}
        \item \textbf{Buffer pool}: area de RAM para cache de datos e indices.
        \item \textit{Cache hit} = dato en RAM (ns); \textit{cache miss} = acceso a disco (ms).
    \end{itemize}
\end{frame}

% ============================================================================
% SECCION 8: COMPLEJIDAD COMPUTACIONAL DE BUSQUEDA
% ============================================================================
\section{Complejidad Computacional de Busqueda}

\begin{frame}{Metodos de acceso}
    \begin{columns}[t]
        \column{0.5\textwidth}
        \textbf{Table Scan}
        \begin{itemize}
            \item O(n): recorre toda la tabla.
            \item Sin indice, o tabla muy pequena.
        \end{itemize}
        \textbf{Indice B-Tree}
        \begin{itemize}
            \item O(log n) para igualdad/rango.
            \item Comun en SQL.
        \end{itemize}
        \column{0.5\textwidth}
        \textbf{Indice Hash}
        \begin{itemize}
            \item O(1) promedio para igualdad exacta.
            \item Ideal para cache (Redis).
        \end{itemize}
    \end{columns}
    \vspace{0.3cm}
    \begin{alertblock}{Nota}
        Los indices aceleran lecturas pero ralentizan escrituras (actualizacion del indice).
    \end{alertblock}
\end{frame}

% ============================================================================
% SECCION 9: ESCALA DE DATOS (UNIDADES)
% ============================================================================
\section{Escala de Datos}

\begin{frame}{De KB a PB}
    \begin{table}
        \centering
        \begin{tabular}{l|c|l}
            \toprule
            Unidad & Bytes & Ejemplo \\
            \midrule
            KB & $10^3$ & Documento texto \\
            MB & $10^6$ & Foto alta resolucion \\
            GB & $10^9$ & Pelicula HD \\
            TB & $10^{12}$ & Disco duro \\
            PB & $10^{15}$ & Grandes empresas (banca, redes sociales) \\
            \bottomrule
        \end{tabular}
    \end{table}
    \begin{itemize}
        \item Facebook ingiere $\approx$ 4 PB diarios.
    \end{itemize}
\end{frame}

% ============================================================================
% SECCION 10: ROLES EN LA GESTION DE DATOS
% ============================================================================
\section{Roles en la Gestion de Datos}

\begin{frame}{Profesionales del dato}
    \begin{description}
        \item[DBA (Database Administrator)]: Operativa diaria, instalacion, parches, backups, tuning.
        \item[Data Engineer]: Pipelines, ETL, infraestructura de datos (Spark, Airflow).
        \item[Data Scientist]: Modelos predictivos y algoritmos de ML.
        \item[Data Analyst]: Informes y visualizaciones.
        \item[Data Steward]: Gobernanza, calidad, cumplimiento.
    \end{description}
\end{frame}

% ============================================================================
% SECCION 11: GOBERNANZA Y ETICA
% ============================================================================
\section{Gobernanza y Etica}

\begin{frame}{Gobernanza de datos}
    \begin{itemize}
        \item Politicas, procesos y estandares para gestion efectiva y segura.
        \item Componentes: catalogo de datos, metadata, data lineage.
        \item Ayuda a prevenir sesgos en IA al rastrear origen de datos.
    \end{itemize}
    \begin{block}{Privacidad (GDPR)}
        \begin{itemize}
            \item Consentimiento explicito.
            \item Derecho al olvido.
            \item Minimizacion y anonimizacion.
        \end{itemize}
    \end{block}
\end{frame}

% ============================================================================
% SECCION 12: IA Y BASES DE DATOS
% ============================================================================
\section{IA y Bases de Datos}

\begin{frame}{Cuando la IA es lenta}
    \begin{exampleblock}{Caso: asistente bancario}
        \begin{enumerate}
            \item La IA solicita datos del usuario.
            \item Tabla sin indice, millones de filas en HDD $\Rightarrow$ consulta de 2 segundos.
            \item El usuario abandona por lentitud.
        \end{enumerate}
    \end{exampleblock}
    \begin{itemize}
        \item Un modelo avanzado es tan rapido como los datos que consume.
        \item Optimizacion de acceso (indices, cache) es critica.
    \end{itemize}
\end{frame}

\begin{frame}{Ventana de contexto en LLMs}
    \begin{itemize}
        \item Memoria a corto plazo del modelo (tokens).
        \item Ej: GPT-3: 4K tokens; Gemini 1.5 Pro: 1M tokens.
        \item No cabe una base de datos completa.
        \item Costo computacional O(n$^2$) para atencion.
    \end{itemize}
\end{frame}

\begin{frame}{Paradigma RAG (Retrieval-Augmented Generation)}
    \begin{enumerate}
        \item Pregunta $\rightarrow$ embedding.
        \item Busqueda en base de datos vectorial (indices ANN: HNSW, IVF).
        \item Fragmentos relevantes + pregunta inyectados en el prompt.
        \item LLM genera respuesta fundamentada.
    \end{enumerate}
    \begin{itemize}
        \item Indices aproximados (ANN) permiten busqueda sublineal en millones de vectores.
    \end{itemize}
\end{frame}

\begin{frame}{Alucinaciones y su relacion con BBDD}
    \begin{block}{Causas desde la base de datos}
        \begin{itemize}
            \item Datos incorrectos o desactualizados.
            \item Recuperacion fallida (indice mal calibrado, embedding pobre).
            \item Timeout en consulta SQL $\Rightarrow$ LLM inventa.
        \end{itemize}
    \end{block}
    \begin{alertblock}{Consecuencia}
        La integridad, actualizacion y velocidad de la BD son primera linea de defensa contra alucinaciones.
    \end{alertblock}
\end{frame}

\begin{frame}{Lecciones para el ingeniero de datos}
    \begin{itemize}
        \item Conoce los indices: B-Tree para SQL, HNSW para vectores.
        \item Mide la latencia: 10 ms vs 2 s decide exito o fracaso.
        \item Monitoriza calidad de datos: linaje y validaciones.
        \item Disena para el peor caso: reintentos, respuestas degradadas.
    \end{itemize}
\end{frame}

% ============================================================================
% SECCION 13: EJERCICIOS PROPUESTOS (DEL SOLUCIONARIO)
% ============================================================================
\section{Ejercicios Propuestos}

\begin{frame}{Ejercicios}
    \begin{enumerate}
        \item Diferencia entre dato, informacion y conocimiento (ejemplo bancario).
        \item Tipos de datos segun estructura y su implicacion en IA.
        \item Ventajas de UTF-8 sobre ASCII/Latin-1.
        \item Comparar CSV, JSON, Avro, Parquet. Elegir para data lake.
        \item Tipos de metadata y su relacion con data lineage.
        \item Funcion de logs y tablas temporales.
        \item Jerarquia de memoria y latencias.
        \item Buffer pool y cache hit.
        \item Complejidad de table scan vs indices B-Tree y hash.
        \item Unidades de almacenamiento (KB a PB) y ejemplo real.
        \item Roles: DBA, Data Engineer, Data Scientist.
        \item Gobernanza y GDPR.
        \item Impacto de latencia en IA (caso asistente).
        \item Ventana de contexto y costo computacional.
        \item Explicar RAG y el papel de indices vectoriales.
        \item Causas de alucinaciones desde BD.
        \item Indices ANN (HNSW, IVF) y ventaja.
        \item Diseno de un motor de recomendacion en tiempo real.
    \end{enumerate}
\end{frame}

% ============================================================================
% SECCION 14: CASO PRACTICO INTEGRADOR
% ============================================================================
\section{Caso Practico: Motor de Recomendacion en Tiempo Real}

\begin{frame}{Diseno de solucion}
    \begin{itemize}
        \item \textbf{Almacenamiento para perfiles activos}: RAM (con SSD como respaldo) para latencia ns.
        \item \textbf{Busqueda por categoria y precio}: Indice B-Tree compuesto sobre (categoria, precio).
        \item \textbf{Busqueda semantica}: Base de datos vectorial (pgvector, Pinecone) con embeddings de productos.
    \end{itemize}
\end{frame}

% ============================================================================
% SECCION 15: RESUMEN
% ============================================================================
\section{Resumen}

\begin{frame}{Lo que hemos aprendido}
    \begin{exampleblock}{Conceptos clave}
        \begin{itemize}
            \item Dato $\rightarrow$ Informacion $\rightarrow$ Conocimiento.
            \item Datos estructurados, semiestructurados, no estructurados.
            \item Codificacion UTF-8, formatos de archivo (Parquet para analitica).
            \item Metadata y data lineage.
            \item Logs y datos temporales.
            \item Jerarquia de memoria: RAM vs disco.
            \item Complejidad de busqueda: Table Scan, B-Tree, Hash.
            \item Escala: KB, MB, GB, TB, PB.
            \item Roles: DBA, Data Engineer, Data Scientist.
            \item Gobernanza, GDPR y sesgo.
            \item IA: ventana de contexto, RAG, indices vectoriales, alucinaciones.
        \end{itemize}
    \end{exampleblock}
\end{frame}

% --- Diapositiva final ---
\begin{frame}
    \centering
    \Huge \textbf{Gracias!}

\end{frame}

\end{document}